\chapter{Graphical expanders from homogeneous cones}\label{sec:examples-homogeneous-expanders}
The one-dimensional corner construction has a counterpart in every dimension. Global existence for Lipschitz graphs is the external input from \cite{EH89}, Theorem 4.6. Scaling and the bounded-slope comparison theorem then give the expander and its uniqueness. This argument concerns entire graphical hypersurfaces; it makes no assertion about the number of nongraphical expanders asymptotic to the same geometric cone.

\begin{theorem}[A unique bounded-slope graphical expander for a Lipschitz cone]\label{thm:lipschitz-cone-graphical-expander}
Let $f:\mathbb R^n\to\mathbb R$, $n\geq1$, be globally $L$-Lipschitz and positively homogeneous of degree one: $f(\lambda x)=\lambda f(x)$ for every $\lambda>0$. There is a unique smooth entire graph $U$ with bounded slope whose dilation flow $\sqrt t\,U(x/\sqrt t)$ converges locally uniformly to $f$ as $t\downarrow0$ and solves graphical mean curvature flow. It satisfies
\[
 \left(\delta_{ij}-\frac{U_iU_j}{1+|DU|^2}\right)U_{ij}
       =\frac{U-x\cdot DU}{2},\qquad |DU|\leq L,
 \qquad |U(x)-f(x)|\leq L\sqrt{2n}.
\]
Its graph is properly embedded and complete. The existence input is \cite{EH89}, Theorem 4.6; the self-similarity and the stated uniqueness follow from the local comparison argument in this project.
\end{theorem}
\begin{proof}
Use Theorem \ref{thm:ecker-huisken-global-graph-existence} to obtain a smooth positive-time entire flow $u$ with initial trace $f$ and $|Du|\leq L$. For every $\lambda>0$ the parabolic rescaling
\[
 u_\lambda(x,t)=\lambda^{-1}u(\lambda x,\lambda^2t)
\]
solves the same graph equation, has the same slope bound, and has initial trace $\lambda^{-1}f(\lambda x)=f(x)$. On each finite common time interval, Theorem \ref{thm:entire-bounded-slope-comparison} gives $u_\lambda=u$. Taking $U=u(\cdot,1)$ therefore yields $u(x,t)=\sqrt t\,U(x/\sqrt t)$. Differentiate this identity and substitute into the graph equation to obtain the displayed elliptic equation with the positive expander sign. The slope estimate is inherited from the flow, and Theorem \ref{thm:graph-lipschitz-half-holder-time} at time one gives the height estimate.

An entire smooth graph is properly embedded: a compact ambient set bounds the base variable, and the graph is closed. Its induced metric dominates the Euclidean base metric. An intrinsic Cauchy sequence thus converges in the base; smoothness on a compact base neighborhood gives convergence in the graph metric. Hence the graph is complete. Finally, another profile in the stated class has a positive-time dilation flow with bounded slope and the same initial trace. Theorem \ref{thm:entire-bounded-slope-comparison} identifies the flows, hence their time-one profiles. Existence is not being reconstructed from the elliptic equation alone.
\end{proof}

\begin{proposition}[Order, symmetries and blow-down of cone expanders]\label{prop:cone-expander-order-symmetry}
Write $U_f$ for the profile in Theorem \ref{thm:lipschitz-cone-graphical-expander}. If $f\leq g$, then $U_f\leq U_g$. For $Q\in O(n)$ one has $U_{f\circ Q}=U_f\circ Q$, and $U_{-f}=-U_f$. In particular all orthogonal symmetries of the cone are inherited. The quantitative blow-down estimate is
\[
 \sup_{x\in\mathbb R^n}|R^{-1}U_f(Rx)-f(x)|
       \leq L\sqrt{2n}/R\qquad(R>0).
\]
If $f(x)=p\cdot x$, then $U_f(x)=p\cdot x$.
\end{proposition}
\begin{proof}
Order follows by applying Theorem \ref{thm:entire-bounded-slope-comparison} to the flows from the two cones. Orthogonal changes of the base and reflection in the height coordinate preserve the graph equation, and give the indicated initial traces. Uniqueness identifies the transformed flows with the corresponding constructed ones. Divide the uniform height estimate in Theorem \ref{thm:lipschitz-cone-graphical-expander} at $Rx$ by $R$ and use homogeneity to prove the blow-down estimate. Finally the linear function is a stationary graphical flow and solves the expander equation because both its Hessian and $U-x\cdot DU$ vanish. Uniqueness gives the last assertion. A vertical translate of a linear profile would not satisfy this origin-centered expander equation unless the added constant were zero.
\end{proof}

\begin{example}[Rotational expanders from circular graphical cones]\label{ex:radial-lipschitz-cone-expander}
For $\ell>0$, the cone $f(x)=\ell|x|$ produces a radial profile $U(x)=w(r)$, $r=|x|$. For $r>0$,
\[
 \frac{w''}{1+(w')^2}+\frac{n-1}{r}w'
       =\frac{w-rw'}2.
\]
At its smooth center, $w'(0)=0$ and, writing $c=w(0)$,
\[
 0<c\leq\ell\sqrt{2n},\qquad w''(0)=\frac{c}{2n},
 \qquad |A|(0)=\frac{c}{2\sqrt n}.
\]
The dilation flow has tip height $c\sqrt t$ and curvature $c/(2\sqrt{nt})$ at its tip. Thus the one-half time scale occurs in radial hypersurface examples in every dimension.
\end{example}
\begin{proof}
Orthogonal symmetry in Proposition \ref{prop:cone-expander-order-symmetry} makes $U$ radial. At $r>0$ its Hessian has eigenvalues $w''$ in the radial direction and $w'/r$ in the $n-1$ orthogonal directions, while the coefficient matrix has radial eigenvalue $(1+(w')^2)^{-1}$ and tangential eigenvalues one. This gives the equation. Smooth radial symmetry gives $DU(0)=0$ and $D^2U(0)=\alpha I$. The equation at the origin gives $n\alpha=c/2$.

Comparison with the zero cone gives $u\geq0$. Its initial trace is positive away from the origin. At any fixed positive time the solution is strictly positive by the classical local strong maximum principle: on positive-time cylinders the graph equation is a smooth uniformly parabolic scalar equation, and an interior zero would force zero at earlier positive times in the connected domain. Chaining overlapping bounded cylinders contradicts locally uniform convergence to the nonzero initial trace. This standard parabolic principle is used as background, not proved here. Hence $c>0$. The height bound at the origin gives $c\leq\ell\sqrt{2n}$. Since the graph metric at the origin is Euclidean, the second fundamental form has $n$ equal principal curvatures in absolute value $c/(2n)$; its norm is $c/(2\sqrt n)$. Dilation by $\sqrt t$ proves both time formulas. No global convexity or classification of all rotational immersions is needed for these conclusions.
\end{proof}

\begin{example}[A nonconvex cone with an expander of changing height sign]\label{ex:saddle-cone-graphical-expander}
In dimension two take $f(x_1,x_2)=|x_1|-|x_2|$. Its unique bounded-slope expander is even in each coordinate and satisfies
\[
 U(x_2,x_1)=-U(x_1,x_2).
\]
It vanishes on both diagonals, is not affine, and takes both positive and negative values. At the origin $U=0$, $DU=0$, $H=0$ and $D^2U=\operatorname{diag}(a,-a)$ for some real $a$. No value or nonvanishing assertion for $a$ is required.
\end{example}
\begin{proof}
The cone is positively one-homogeneous and $\sqrt2$-Lipschitz. Its separate even symmetries and sign reversal under swapping coordinates transfer to $U$ by Proposition \ref{prop:cone-expander-order-symmetry}. They force the stated vanishing on the diagonals and the values of $U$, its gradient, and the shape of its Hessian at the origin. The expander equation gives zero mean curvature there. The uniform height estimate gives $|U-f|\leq2\sqrt2$, so $U(R,0)>0$ and $U(0,R)<0$ for large $R$. If $U$ were affine, its blow-down would be a linear function, contradicting the explicit nonaffine cone. Moreover its scalar mean curvature with the upward orientation satisfies $H(x_2,x_1)=-H(x_1,x_2)$, by the formula $H=-\operatorname{div}(DU/v)$. If $H$ had one sign it would vanish identically. The expander equation would then give $U=x\cdot DU$, making $U$ homogeneous of degree one; smoothness at the origin forces such a function to be linear, a contradiction. Thus the mean curvature changes sign. This example does not classify its zero set beyond the two diagonals, or infer a sign for each principal curvature away from the origin.
\end{proof}

\begin{corollary}[Convex cones have strictly positive expander defect]\label{cor:convex-cone-positive-expander-defect}
If the cone $f$ in Theorem \ref{thm:lipschitz-cone-graphical-expander} is convex and not linear, then its profile satisfies
\[
 U_f(x)>f(x),\qquad U_f(x)-x\cdot DU_f(x)>0,\qquad H<0
\]
for the upward orientation. In particular $U_f(0)>0$. These conclusions hold in every dimension and distinguish convex cones from Example \ref{ex:saddle-cone-graphical-expander}.
\end{corollary}
\begin{proof}
Proposition \ref{prop:convex-graph-data-upward-motion} gives strict height increase and $u_t>0$ for the constructed flow. Its self-similar expression gives $2u_t(x,1)=U_f(x)-x\cdot DU_f(x)$. This proves the defect inequality and the sign of $H=-u_t/v$. A positively one-homogeneous continuous function has $f(0)=0$, so the strict height inequality yields $U_f(0)>0$. The proof does not require a separate convexity-preservation theorem for the flow.
\end{proof}
